\originalsection{多频运动与准周期}
一维周期轨道只有一个作用量. 若$n$自由度系统的HJ特征函数可完全分离为$W=\sum_iW_i(q_i,\alpha)$, 每个分离运动都是振荡或转动, 可沿相应的基本闭合回路定义$I_i=(2\pi)^{-1}\oint p_i\dd q_i$. 这里的回路只改变第$i$个分离坐标, 并不要求真实的多维轨道已经闭合. 假设作用量对分离常数的Jacobi矩阵可逆, 才能局部把$\alpha$写成$I$的函数, 并得到$H=H(I_1,\ldots,I_n)$.

令$\vartheta_i=\partial W/\partial I_i$. 沿第$j$个基本回路,
\[
\Delta_j\vartheta_i
=\frac{\partial}{\partial I_i}\oint_j p_j\dd q_j
=2\pi\delta_{ij}.
\]
因此$\vartheta_i$是模$2\pi$的角坐标, 方程成为$\dot I_i=0$, $\dot\vartheta_i=\omega_i(I)=\partial H/\partial I_i$. 固定$I$后, 所有角构成$n$维环面; 运动是在环面上以恒定角速度前进. 本册用角频率$\omega_i$, 原稿用未除$2\pi$的作用量及每周期增1的角, 相应频率为$\nu_i=\omega_i/(2\pi)$.

完整角状态周期的条件是存在$T>0$使所有$\omega_iT\in2\pi\mathbb Z$. 对两非零频率, 这等价于$\omega_1/\omega_2$为有理数. 例如二维各向异性谐振子有$x=A\cos(\omega_1t+\delta_1)$, $y=B\cos(\omega_2t+\delta_2)$; 若比值为$p/q$且$p,q$互素, 可取$T=2\pi p/\omega_1=2\pi q/\omega_2$. 比值无理时, 两个坐标各自周期, 完整相状态却没有共同周期, 称准周期. 对非退化振幅, 后者的角轨道在二维环面上稠密. 多频情况下还须区分各频率都可公度与仅存在某些整数共振关系: 一条共振关系可以把运动限制到较低维子环面, 不必使整条轨道闭合.

\originalsection{三维Kepler的分离与角作用量}
前面的平面解利用了角动量预先确定轨道平面. 若不预先选择这个平面, 在球坐标$(r,\theta,\phi)$中,
\[
H=\frac1{2\mu}\left(p_r^2+\frac{p_\theta^2}{r^2}
+\frac{p_\phi^2}{r^2\sin^2\theta}\right)-\frac\alpha r.
\]
设$W=W_r(r)+W_\theta(\theta)+J_z\phi$, 循环坐标$\phi$给$p_\phi=J_z$. 将HJ方程乘$2\mu r^2$, 可把角部分与径向部分分开, 因而引入常数$J^2$使
\[
p_\theta^2+\frac{J_z^2}{\sin^2\theta}=J^2,
\qquad p_r^2=2\mu E+\frac{2\mu\alpha}{r}-\frac{J^2}{r^2}.
\]
这两个常数并非另造的积分: $J$是角动量大小, $J_z$是其$z$分量. 因而$|J_z|\le J$.

为把符号和回路方向固定下来, 先在$0<J_z<J$的局部坐标片上计算, 且取非圆有界轨道$E<0$, $0<e<1$. 极角转向点满足$\sin\theta_0=c=J_z/J$, 另一点为$\pi-\theta_0$. 转动作用量为$I_\phi=J_z$, 极角振荡作用量为
\[
I_\theta=\frac{2J}{\pi}\int_{\theta_0}^{\pi/2}
\sqrt{1-\frac{c^2}{\sin^2\theta}}\dd\theta.
\]
设$b=\sqrt{1-c^2}$, 令$\cos\theta=b\sin\psi$. 原积分从$\theta_0$到$\pi/2$对应$\psi$从$\pi/2$到0, 反转界限后有
\[
\int_{\theta_0}^{\pi/2}\sqrt{1-\frac{c^2}{\sin^2\theta}}\dd\theta
=\int_0^{\pi/2}\frac{b^2\cos^2\psi}{1-b^2\sin^2\psi}\dd\psi
=\frac\pi2-c^2\int_0^{\pi/2}\frac{\dd\psi}{1-b^2\sin^2\psi}.
\]
最后一积分用$u=\tan\psi$变成$\int_0^\infty\dd u/(1+c^2u^2)=\pi/(2c)$. 所以$I_\theta=J-J_z$, $I_\theta+I_\phi=J$. 若允许$J_z<0$而按运动方向定义非负转动作用量, 应改写为$I_\phi=|J_z|$, $I_\theta=J-|J_z|$, 角变换也要调整符号; 不能把上述局部公式无条件用于逆行轨道.

径向作用量已在上一节算出$I_r=\sqrt{\mu\alpha a}-J$. 因而三维Kepler哈密顿量为
\[
H=-\frac{\mu\alpha^2}{2\Lambda^2},\qquad
\Lambda=I_r+I_\theta+I_\phi=\sqrt{\mu\alpha a}.
\]
在上述坐标片上三个角频率都等于$\mu\alpha^2/\Lambda^3$. 对一般中心势, $H$只依赖$I_r$和$J=I_\theta+I_\phi$, 故两角频率相同, 径向频率却不必与它们相同. Kepler多出这一次简并, 对应非圆有界轨道的闭合.

\originalsection{轨道的五个常数与运动相位}
作用量形式把椭圆形状和方位分得很清楚. 定义新动量$K_1=I_\phi$, $K_2=I_\theta+I_\phi=J$, $K_3=\Lambda$, 并用第二型生成函数
\[
F_2=(\vartheta_\phi-\vartheta_\theta)K_1
+(\vartheta_\theta-\vartheta_r)K_2+\vartheta_rK_3.
\]
求旧动量可核对$I_\phi=K_1$, $I_\theta=K_2-K_1$, $I_r=K_3-K_2$; 新坐标为$\chi_1=\vartheta_\phi-\vartheta_\theta$, $\chi_2=\vartheta_\theta-\vartheta_r$, $\chi_3=\vartheta_r$. 由于$H$只依赖$K_3$, $\dot\chi_1=\dot\chi_2=0$, 只有$\chi_3$随时间前进.

要把常角认成空间方位角, 还须固定特征函数积分的参考点. 令$W_\theta$从升交点的$\theta=\pi/2$起积分, 在每一转向点沿实际$p_\theta$分支连续延拓; $W_r$从近日点$r_-$起积分. 记轨道平面内从升交点量起的角为$u$, 真近点角为$f$, 则$u=\omega+f$. 球坐标的几何关系及角动量给
\[
\cos\theta=\sin i\sin u,\qquad
\frac{\dd\phi}{\dd u}=\frac{J_z}{J\sin^2\theta},\qquad
p_\theta=J\frac{\dd\theta}{\dd u}.
\]
在固定另一角常量时对$p_\theta$求偏导, 分别为$J/p_\theta$和$-J_z/(p_\theta\sin^2\theta)$. 由上式积分得$\partial_JW_\theta=u$与$\partial_{J_z}W_\theta=-(\phi-\Omega)$. 固定$\Lambda$时$E$固定, 径向偏导为$\partial_Jp_r=-J/(r^2p_r)$; 又沿轨道有$\dd f=J\dd r/(r^2p_r)$, 故$\partial_JW_r|_\Lambda=-f$. 下限$r_-$虽然依赖$J$, 但端点$p_r(r_-)=0$, 不增加端点项. 各式先在非转向的局部分支成立, 再沿轨道连续延拓. 因而$\chi_1=\phi+\partial_{J_z}W_\theta=\Omega$, $\chi_2=\partial_J(W_\theta+W_r)=u-f=\omega$, 不是只凭它们守恒便作认定. 三个常作用量及两个常角分别给
\[
\cos i=\frac{K_1}{K_2},\qquad
a=\frac{K_3^2}{\mu\alpha},\qquad
e=\sqrt{1-\frac{K_2^2}{K_3^2}},\qquad
\Omega=\chi_1,\qquad \omega=\chi_2.
\]
$i$是轨道平面对参考平面的倾角, $\Omega$是升交点经度, $\omega$是轨道平面内从升交点指向近日点的角. 两个方位角确定平面, $a,e$确定大小和形状, $\omega$确定椭圆在平面内的朝向. 它们共五个常数, 确定的是几何轨道. 要确定质点何时在轨道哪个位置, 还需第六个积分常数, 即$\chi_3$的初始相位或近日点通过时刻. 把“五个轨道参数”说成“全部初始条件”会漏掉这一点. 当$e=0$时近日点角失去意义, 当$i=0$时升交点也失去意义; 这是这些参数的坐标退化, 不是圆轨道或平面轨道不能研究.

这些常数也可从守恒矢量看出. 定义偏心率矢量$\vect e=(\vect p\times\vect J)/(\mu\alpha)-\hat{\vect r}$, 即归一化的Laplace--Runge--Lenz矢量. 由于$\dot{\vect p}=-\alpha\hat{\vect r}/r^2$且$\dot{\vect J}=0$,
\[
\frac{\dot{\vect p}\times\vect J}{\mu\alpha}
=\frac{\vect p-\hat{\vect r}(\hat{\vect r}\cdot\vect p)}{\mu r}
=\frac{\dd\hat{\vect r}}{\dd t},
\]
故$\dot{\vect e}=0$. 它满足$\vect e\cdot\vect J=0$, 而不是叉积为零; 8.223第8讲第8页在此有笔误. 又由$\vect e\cdot\hat{\vect r}=J^2/(\mu\alpha r)-1$重得椭圆轨道方程, $|\vect e|=e$, 方向指向近日点. $\vect J$给平面法向, $\vect e$给平面内的朝向. 这些矢量受$\vect e\cdot\vect J=0$及$e^2=1+2EJ^2/(\mu\alpha^2)$约束, 不能把每个分量都另算成独立守恒量.
